\section{Experiments}

\subsection{Compared Systems}

We evaluated six logged primary entries: an LSTM baseline, a vanilla Transformer, a Transformer without constraints, the proposed neural-symbolic rule-guided Transformer, a masked-infilling task model, and a soprano-to-SATB task model. Each training run wrote checkpoints, evaluation summaries, MusicXML examples, and consolidated metric rows. These artifacts are stored in the repository-level run, result, and generated-score directories.

\subsection{Ablation Design}

Four logged ablations were used to test whether the proposed system's behaviour came from specific neural-symbolic components rather than from model size alone. The no-harmony ablation removed automatic harmonic conditioning. The no-refinement ablation replaced iterative masked rewriting with a single prediction pass. The no-rule ablation disabled symbolic rule-guided decoding while retaining the neural architecture. The no-voice-relation ablation removed the voice-relation attention component. All ablation rows were generated from local runs and logged through the same evaluation scripts as the primary systems.

\subsection{Artifact Ingestion}

Tables and figures were generated only from stored CSV and JSON files. The inspected metric file contains one CPU smoke row, five non-smoke primary rows, and four non-smoke ablation rows used in the current ablation table, plus historical exploratory rows retained in the CSV audit trail. The smoke row is retained as a software check and is not used as a final experimental result. The logged vanilla Transformer and Transformer-without-constraints rows point to the same no-constraints checkpoint, so they are treated as the same baseline when interpreting ablations.

\subsection{Generated Scores and Expert Materials}

The primary run generated score-level MusicXML examples for generated and ground-truth SATB pairs. A separate expert-evaluation package was prepared with anonymized score PDFs, paired A/B comparison PDFs, and blank rating forms. The expert package is intended for score reading, not audio listening; playback may be used only as a notation aid. At the time of this manuscript update, no completed expert ratings were present, so the paper reports the expert-evaluation design but does not report expert scores.
